\documentclass[10pt]{article}
\usepackage[T1]{fontenc}
\usepackage[utf8]{inputenc}
\usepackage[french]{babel}
\usepackage{amsmath,amssymb,mathrsfs,mathpazo}
\usepackage{geometry}
\geometry{a4paper,margin=2cm}
\usepackage{graphicx,array,tabularx,tabularray,colortbl,multirow,multicol,booktabs,enumitem,pifont,ulem,fancybox,diagbox,cancel}
\usepackage{pstricks,pstricks-add,pst-plot,pst-tree,pst-node,pst-text,pst-func,pst-math,pst-bspline,pst-3dplot,pst-eucl,pst-solides3d}
\let\pscalculate\undefined
\usepackage{pst-bezier}
\usepackage{tikz,pgfplots,tkz-tab}
\usetikzlibrary{arrows,arrows.meta,patterns,calc,positioning,trees,decorations.markings,shapes.geometric,shapes.misc,shapes.arrows,intersections,decorations.pathmorphing,decorations.pathreplacing,backgrounds,shadows,babel}
\pgfplotsset{compat=1.18}
\usepgfplotslibrary{fillbetween}
\usepackage[np]{numprint}
\usepackage{listings,makecell,pas-tableur,twoopt}
\let\barre\undefined
\usepackage[pstricks]{tablvar}
\definecolor{darkgreen}{RGB}{0,100,0}
\providecommand{\red}{\color{red}}
\providecommand{\blue}{\color{blue}}
\providecommand{\green}{\color{green}}
\providecommand{\black}{\color{black}}
\providecommand{\white}{\color{white}}
\providecommand{\gray}{\color{gray}}
\providecommand{\cyan}{\color{cyan}}
\providecommand{\magenta}{\color{magenta}}
\providecommand{\yellow}{\color{yellow}}
\providecommand{\violet}{\color{violet}}
\providecommand{\orange}{\color{orange}}
\providecommand{\euro}{\texteuro}
\providecommand{\EUR}{\texteuro}
\providecommand{\ieme}{\textsuperscript{e}}
\providecommand{\degres}{\ensuremath{^\circ}}
\providecommand{\degree}{\ensuremath{^\circ}}
\providecommand{\vv}[1]{\overrightarrow{#1}}
\providecommand{\Coord}[1]{#1}
\providecommand{\decosix}{}
\providecommand{\starredbullet}{\ensuremath{\bullet}}
\pagestyle{empty}
\setlength{\parindent}{0pt}
\usepackage{adjustbox}\DeclareUnicodeCharacter{2236}{:}\DeclareUnicodeCharacter{2190}{\ensuremath{\leftarrow}}
\providecommand{\R}{\mathbb{R}}
\providecommand{\N}{\mathbb{N}}
\providecommand{\Z}{\mathbb{Z}}
\providecommand{\Q}{\mathbb{Q}}
\providecommand{\D}{\mathbb{D}}
\providecommand{\C}{\mathbb{C}}
\providecommand{\e}{{\mathrm{e}}}
\providecommand{\vect}[1]{\overrightarrow{#1}}
\providecommand{\vectt}[1]{\overrightarrow{\mathrm{#1}}}
\providecommand{\vvt}[1]{\overrightarrow{\mathrm{#1}}}
\providecommand{\np}[2][]{{#2~#1}}
\providecommand{\vv}[1]{{\overrightarrow{#1}}}
\renewcommand{\d}{{\mathrm{d}}}
\providecommand{\degres}{{^\circ}}
\providecommand{\degre}{{^\circ}}
\providecommand{\blue}{\color{blue}}
\providecommand{\red}{\color{red}}
\providecommand{\enskip}{\quad}
\providecommand{\llbracket}{\lbrack\!\lbrack}
\providecommand{\rrbracket}{\rbrack\!\rbrack}
\providecommand{\euro}{{\text{€}}}
\providecommand{\no}{{\text{n°}}}
\providecommand{\interval}[2]{{[#1;#2]}}
\providecommand{\intervFO}[2]{{[#1;#2[}}
\providecommand{\dotfill}{\ldots}

\providecommand{\ds}{\displaystyle}
\providecommand{\textding}[1]{#1}
\providecommand{\Oij}{\(\mathrm{O};\vec{\imath},\vec{\jmath}\)}
\providecommand{\Oijk}{\(\mathrm{O};\vec{\imath},\vec{\jmath},\vec{k}\)}
\providecommand{\Ouv}{\(\mathrm{O};\vec{u},\vec{v}\)}
\providecommand{\intervFF}[2]{\ensuremath{[#1;#2]}}
\renewcommand{\euro}{\ifmmode\text{\texteuro}\else\texteuro\fi}
\DeclareUnicodeCharacter{00B0}{\ifmmode^{\circ}\else\textdegree\fi}
\DeclareUnicodeCharacter{20AC}{\ifmmode\text{\texteuro}\else\texteuro\fi}
% Énoncés : Éducation nationale. Transcription et corrigés : APMEP, Denis Vergès et contributeurs. Découpage et classement : CoopMaths, Sébastien Lozano, Éric Elter et contributeurs. Adaptation de la présentation : Maths Explorer.
% CC BY-SA : https://coopmaths.fr/www/about/
% https://coopmaths.fr/alea/static/e3c/2020/tex/e3c_2020_00_sujet41_4_cor.tex
\begin{document}
\noindent\textit{APMEP / CoopMaths. CC BY-SA. Présentation adaptée par Maths Explorer.}\par\medskip
\( f(x) = \dfrac{x^2 - 4x + 8}{x - 2} \).

\paragraph{1.} Sur l'intervalle \( ]-\infty \,;\, 2[ \), \( f(x) = 0 \iff x^2 - 4x + 8 = 0 \).
\[
\Delta = 16 - 4 \times 8 = 16 - 32 = -16 < 0.
\]
L'équation n'a pas de solutions réelles dans \( ]-\infty \,;\, 2[ \).

\paragraph{2.}

\paragraph{a.} Le dénominateur étant non nul, puisque \( x < 2 \), la fonction quotient est dérivable et :
\begin{align*}
f'(x) &= \dfrac{(2x - 4)(x - 2) - (x^2 - 4x + 8)}{(x - 2)^2} \\
&= \dfrac{2x^2 - 4x - 4x + 8 - x^2 + 4x - 8}{(x - 2)^2} \\
&= \dfrac{x^2 - 4x}{(x - 2)^2}.
\end{align*}

\paragraph{b.} Le dénominateur étant positif, le signe de \( f'(x) \) est celui de \( x^2 - 4x = x(x - 4) \).

Ce trinôme a pour racines 0 et 4 et on sait que ce trinôme est positif, sauf sur \( ]0 \,;\, 4[ \).

Dans notre cas \( f'(x) \) est positif sur \( ]-\infty \,;\, 0[ \) et négatif sur \( ]0 \,;\, 2[ \).

La fonction \( f \) est donc croissante sur \( ]-\infty \,;\, 0[ \) et décroissante sur \( ]0 \,;\, 2[ \).

\paragraph{3.} On sait qu'une équation de \( \mathcal{D} \) est :
\[
y - f(1) = f'(1)(x - 1).
\]

Avec \( f(1) = \dfrac{1 - 4 + 8}{1 - 2} = - 5 \) et \( f'(1) = \dfrac{1 - 4}{(-1)^2} = -3 \), l'équation devient :
\begin{align*}
y + 5 &= -3(x - 1) \\
y &= -3x + 3 - 5 \\
y &= -3x - 2.
\end{align*}

\paragraph{4.}

\begin{center}
\psset{unit=0.65cm}
\begin{pspicture}(-10,-10)(10,3)
\psaxes[linewidth=1.25pt,Dx=2,Dy=2]{->}(0,0)(-10,-10)(10,3)
\def\Func{x 2 exp 4 x mul sub 8 add x 2 sub div}
\psplot[plotpoints=2000,linewidth=0.85pt,linecolor=red]{-7.583}{1.583}{\Func}
\def\FuncTwo{-3 x mul 2 sub}
\psplot[plotpoints=2000,linewidth=0.85pt,linecolor=blue]{-1.667}{2.667}{\FuncTwo}
\uput[u](-2,2){$\mathscr{D}$}
\end{pspicture}
\end{center}

\end{document}
